\documentclass[10pt]{article}

\usepackage[margin=0.8in]{geometry}
\usepackage{amsmath,amssymb,amsthm}
\usepackage[T1]{fontenc}
\usepackage{lmodern}

\newtheorem{lemma}{Lemma}

\begin{document}

\begin{center}
    {\large\bfseries Alternative proof of Lemma 3}
\end{center}

\medskip

Recall that
\[
T_1(x)=x_1+x_2,
\qquad
\Pi=\operatorname{cone}\{(1,0),(1,-1)\},
\qquad
\widehat A=\operatorname{cl}(A-\Pi),
\]
and define
\[
\mathcal F=\{a\in\mathbf R^2:0\le a_2\le a_1\}.
\]

\setcounter{lemma}{2}
\begin{lemma}[Calibration of the competitor]
Suppose that $A\subseteq\mathbf R^2$ is closed, convex, and coordinatewise lower,
$P_0(A)=1-\alpha$, and $A$ has a finite-support direction $a$ satisfying
$a_2>a_1$. Then there exists a finite number $b$ such that
\[
P_0\bigl(\widehat A\cap\{x:T_1(x)\le b\}\bigr)=1-\alpha.
\]
For any such $b$, define
\[
A'=\widehat A\cap\{x:T_1(x)\le b\}.
\]
Then $A'$ is closed, convex, and coordinatewise lower, satisfies
$P_0(A')=1-\alpha$, and every nonzero finite-support direction of $A'$ lies in
$\mathcal F$. Consequently, the test $\phi_{A'}(x)=I\{x\notin A'\}$ has level
$\alpha$, attained at $\mu=0$.
\end{lemma}

\begin{proof}
By Lemma 2,
\[
P_0(\widehat A)>P_0(A)=1-\alpha.
\]
The map
\[
t\longmapsto P_0\bigl(\widehat A\cap\{x:T_1(x)\le t\}\bigr)
\]
is continuous because $P_0\{T_1(X)=t\}=0$ for every $t$. It converges to
zero as $t\to-\infty$ and to $P_0(\widehat A)$ as $t\to\infty$. Hence there
exists a finite $b$ for which the displayed probability equals $1-\alpha$.

The set $A'$ is closed, convex, and coordinatewise lower because it is the
intersection of two sets with these properties. Its null probability equals
$1-\alpha$ by construction.

It remains to characterize its finite-support directions. We first show that
$A'$ is $\Pi$-lower. Let $x\in A'$ and $p\in\Pi$. Lemma 2 implies that
$x-p\in\widehat A$. Moreover, $T_1(p)=p_1+p_2\ge0$ for every $p\in\Pi$, and
therefore
\[
T_1(x-p)=T_1(x)-T_1(p)\le T_1(x)\le b.
\]
Thus $x-p\in A'$.

Now let $a\ne0$ satisfy $h_{A'}(a)<\infty$, and fix $x\in A'$. Since $A'$ is
coordinatewise lower, $x-te_j\in A'$ for every $t\ge0$ and $j=1,2$. If
$a_j<0$, then
\[
a'(x-te_j)=a'x-ta_j\longrightarrow+\infty,
\]
contradicting $h_{A'}(a)<\infty$. Hence $a_1\ge0$ and $a_2\ge0$.

Because $(1,-1)\in\Pi$ and $A'$ is $\Pi$-lower,
\[
x-t(1,-1)\in A'
\qquad\text{for every }t\ge0.
\]
If $a_2>a_1$, then
\[
a'\{x-t(1,-1)\}=a'x+t(a_2-a_1)\longrightarrow+\infty,
\]
again contradicting finite support. Therefore $a_2\le a_1$. We have shown
that every nonzero $a$ with $h_{A'}(a)<\infty$ satisfies
\[
0\le a_2\le a_1,
\]
so $a\in\mathcal F$.

Finally, the complement of $A'$ is coordinatewise upper. For every
$\mu\in\Lambda_0$, we have $\mu\le0$. Coupling $X=\mu+Z$ with
$Z\sim N_2(0,I_2)$ gives
\[
P_\mu(X\notin A')\le P_0(Z\notin A')=\alpha.
\]
Equality holds at the finite null point $\mu=0$, which proves the last claim.
\end{proof}

\end{document}
